\section{Equality at the original germ and global consequences}
\label{sec:conclusion}

We complete the return from the cone to the original singularity.
The lci case of the gap theorem is already known
\cite[Theorem~1.3]{Liu22}. We include the hypersurface argument
needed here, so that the source of strictness in the equality
case is explicit.

\subsection{A weighted hypersurface bound}

\begin{lemma}[A hypersurface test]\label{end:hypersurface-test}
Let \(x\in X\) be a klt hypersurface germ of dimension \(n\),
with equation \(F\) in regular ambient parameters
\(z_0,\ldots,z_n\). For positive rational weights
\(w_0,\ldots,w_n\), write \(d_F\) for the weighted order of \(F\).
Then
\begin{equation}\label{end:eq-hypersurface-test}
 \nvol(x,X)\le
           \frac{d_F\bigl(\sum_{i=0}^n w_i-d_F\bigr)^n}
                {\prod_{i=0}^n w_i}.
\end{equation}
\end{lemma}

\begin{proof}
Both sides of \eqref{end:eq-hypersurface-test} are unchanged by
a common scaling of the weights. We may therefore assume
that they are positive integers. Choose a positive common
multiple \(D_1\) of the \(w_i\), and in \(\cO_{X,x}\) put
\[
                       I=(z_i^{D_1/w_i}:0\le i\le n).
\]
This ideal is \(\fm_x\)-primary.

The ambient weighted valuation has log discrepancy
\(\sum_i w_i\), gives value \(d_F\) to \(F\), and value
\(D_1\) to the ambient ideal defining \(I\). Hypersurface
inversion of adjunction \cite[Theorem~1.1]{EinMustata04}, applied to the smooth ambient
space with \(F\) of coefficient one, gives
\begin{equation}\label{end:eq-threshold-test}
                   \lct_X(I)\le
                   \frac{\sum_i w_i-d_F}{D_1}.
\end{equation}
In particular the numerator on the right is positive,
because \(X\) is klt and \(I\) is primary.

Let \(J_t\) be the quotient weighted-order filtration on
the hypersurface. The elementary monomial floor inequality
gives
\[
                  J_{(k+n+1)D_1}\subset I^k\subset J_{kD_1}.
\]
Indeed, for a monomial of weight at least \((k+n+1)D_1\),
the sum of the integers
\(\lfloor \alpha_i w_i/D_1\rfloor\) is at least \(k\).
The associated graded of the quotient filtration is
\[
 \CC[z_0,\ldots,z_n]/(\operatorname{in}_w F),\qquad
 H(t)=\frac{1-t^{d_F}}{\prod_i(1-t^{w_i})}.
\]
This formula holds even if \(\operatorname{in}_w F\) is
reducible or nonreduced: it is a nonzero equation in a
polynomial ring. The two filtration inclusions and the
leading Hilbert coefficient give
\begin{equation}\label{end:eq-multiplicity}
                       e(I)=\frac{D_1^n d_F}{\prod_i w_i}.
\end{equation}

Choose a divisor \(E\) computing \(\lct_X(I)\), and put
\(q=\ord_E(I)>0\). Its center is \(x\), since \(I\) is
primary. The inclusion \(I^k\subset\fa_{kq}(E)\) implies
\[
             q^n\vol(E)\le e(I).
\]
Testing the normalized-volume infimum with \(E\), then
using \eqref{end:eq-threshold-test} and
\eqref{end:eq-multiplicity}, yields
\[
 \nvol(x,X)\le A_X(E)^n\vol(E)
            \le \lct_X(I)^n e(I),
\]
which is \eqref{end:eq-hypersurface-test}.
\end{proof}

\begin{proposition}[Hypersurface equality]\label{end:hypersurface}
If \(x\in X\) is a singular klt hypersurface germ of
dimension \(n\ge2\), then
\[
                        \nvol(x,X)\le2(n-1)^n.
\]
If equality holds, the quadratic part of a local equation
is nondegenerate, and \((X,x)\) is analytically an ordinary
double point.
\end{proposition}

\begin{proof}
Let \(q=\mult_x(F)\ge2\). Weights all equal to one in
Lemma~\ref{end:hypersurface-test} give
\[
                        \nvol(x,X)\le q(n+1-q)^n,
\]
and klt implies \(q<n+1\). The function
\(b(t)=t(n+1-t)^n\) satisfies
\[
 b'(t)=(n+1)(1-t)(n+1-t)^{n-1}<0
                         \quad(1<t<n+1).
\]
Thus \(q>2\) gives a strict bound below \(2(n-1)^n\).

Suppose \(q=2\) and the quadratic part is degenerate.
After a linear change of coordinates, choose a direction
absent from the quadratic part. Give it weight
\(1-\epsilon\), and give every other direction weight one,
where \(0<\epsilon<1/3\) is rational. Every term of
ordinary degree at least three has weight greater than
two, so \(d_F=2\). The bound becomes
\[
             \frac{2(n-1-\epsilon)^n}{1-\epsilon}.
\]
Its logarithmic derivative at \(\epsilon=0\) is
\(-n/(n-1)+1=-1/(n-1)<0\). It is therefore strictly
smaller than \(2(n-1)^n\) for sufficiently small
\(\epsilon>0\). Equality forces a nondegenerate
quadratic part, and the holomorphic Morse Lemma gives
the asserted analytic germ.
\end{proof}

\subsection{Completion of the induction}

\begin{proof}[Proof of Theorem~\ref{thm:main}]
The cases \(n=2,3,4\) were established in
Section~\ref{sec:cones}, the last by
Theorem~\ref{inp:fourfold}. Suppose the theorem is
known in lower dimensions and let \(n=N+1\ge5\).
The supremum argument of Theorem~\ref{boot:upper}
gives \(M_n\le p_n\), hence the asserted inequality
for every singular germ.

Now suppose \(d(x,X)=p_n\). Apply
Proposition~\ref{cone:reduction} to obtain a
K-polystable cone \(C\) with the same density and
\(\edim(X,x)\le\edim(C,o)\).
The equality argument of
Theorem~\ref{boot:equality-cone} makes \(C\) an
ordinary double point. Therefore
\(\edim(X,x)\le n+1\).

Since \(x\) is singular, its embedding dimension is
exactly \(n+1\), and it is a hypersurface germ.
To see this algebraically, start with an embedding
in a smooth ambient variety and use equations with
independent linear parts to cut the ambient germ
down to dimension \(n+1\). The remaining ideal is a
height-one prime in a regular local ring, so is
principal. Proposition~\ref{end:hypersurface}
shows that its equation has nondegenerate quadratic
part and hence that the analytic germ is ODP.

Conversely, Lemma~\ref{end:odp-value} gives equality
at every analytic ODP. This proves both directions
of the equality statement and completes the induction.
\end{proof}

\subsection{The global comparison}

\begin{proof}[Proof of Corollary~\ref{cor:global}]
For every closed point \(x\) of a K-semistable
\(\QQ\)-Fano \(n\)-fold \(V\), the local-to-global
comparison gives
\[
                    (-K_V)^n
                    \le\left(\frac{n+1}{n}\right)^n
                         \nvol(x,V).
\]
This follows from \cite[Theorem~21]{Liu18};
K-semistability and Ding semistability agree in
this setting by \cite[Theorem~13(2)]{Liu18}.
If \(V\) is singular, apply
Theorem~\ref{thm:main} at a singular closed point:
\[
 (-K_V)^n\le
 \left(\frac{n+1}{n}\right)^n2(n-1)^n
 =2\left(\frac{n^2-1}{n}\right)^n<2n^n.
\]

If \(V\) is smooth and \(V\not\cong\PP^n\), the
theorem of Li--Miao gives \((-K_V)^n\le2n^n\);
its equality cases are a smooth quadric and
\(\PP^1\times\PP^{n-1}\)
\cite[Theorem~1.1]{LM25}.
Both have Kähler--Einstein metrics, given by the homogeneous
quadric metric and the suitably scaled Fubini--Study product metric,
respectively. They are therefore K-polystable
\cite[Theorem~1.1]{Berman16}. Their degrees are
\[
 (-K_{Q^n})^n=(nH)^n=2n^n,\qquad
 (-K_{\PP^1\times\PP^{n-1}})^n
                =(2H_1+nH_2)^n=2n^n.
\]
For \(n=2\), a smooth quadric surface is
\(\PP^1\times\PP^1\). The strict singular bound
excludes every other equality case.
\end{proof}
